\chapter{Complete Python Code}

\begin{lstlisting}[language=Python,
    caption={Complete Python implementation for Batch 9 BPSK system},
    basicstyle=\ttfamily\small,
    breaklines=true,
    showstringspaces=false]
import numpy as np
import matplotlib.pyplot as plt

from scipy.signal import lfilter
from scipy.special import erfc

# ============================================================
# CONFIGURATION
# ============================================================

np.random.seed(42)

A = 1.0
fc = 15e3
Rb = 1e3
fs = 200e3
bits_plot = np.array([1, 0, 1, 1, 0, 0, 1, 0, 1, 0])
h = np.array([0.2, 0.6, 0.2])
n_bits = 10**5
EbN0_dB_range = np.array([2, 4, 6, 8, 10, 12, 14])

# ============================================================
# PART A - BPSK SYMBOL GENERATION
# ============================================================

Tb = 1 / Rb
N = int(fs / Rb)
t = np.arange(0, Tb, 1 / fs)

s1 = A * np.cos(2 * np.pi * fc * t)
s0 = -A * np.cos(2 * np.pi * fc * t)

Eb_num = np.sum(s1**2) / fs
Eb_theory = (A**2 * Tb) / 2

print("Numerical Eb:", Eb_num)
print("Theoretical Eb:", Eb_theory)

# 10-bit waveform generation
waveform = np.array([])

for bit in bits_plot:
    waveform = np.concatenate((waveform, s1 if bit == 1 else s0))

time_axis = np.arange(len(waveform)) / fs

# ============================================================
# PART B - FFT, PSD, AND PARSEVAL
# ============================================================
X = np.fft.fft(waveform)
X_shift = np.fft.fftshift(X)

freq = np.fft.fftfreq(len(waveform), d=1 / fs)
freq_shift = np.fft.fftshift(freq)

PSD = (np.abs(X_shift) ** 2) / len(waveform)

f_lower_null = fc - Rb
f_upper_null = fc + Rb

BW = 2 * Rb

E_time = np.sum(waveform**2) / fs
L = len(waveform)

E_freq = (1 / fs) * (1 / L) * np.sum(np.abs(X) ** 2)
parseval_error = abs((E_time - E_freq) / E_time) * 100

print("Lower first null:", f_lower_null)
print("Upper first null:", f_upper_null)
print("Null-to-null bandwidth:", BW)
print("Parseval error (%):", parseval_error)

# ============================================================
# PART C - CHANNEL + CORRELATOR RECEIVER
# ============================================================
EbN0_dB_test = 6
EbN0_lin_test = 10 ** (EbN0_dB_test / 10)
N0_test = Eb_theory / EbN0_lin_test
sigma2_test = N0_test * fs / (2 * Rb)
sigma_test = np.sqrt(sigma2_test)

tx_filtered = lfilter(h, 1, waveform)
noise = sigma_test * np.random.randn(len(tx_filtered))
rx = tx_filtered + noise

def correlator_receiver(rx_signal, ref_signal, N):
    detected_bits = []
    correlator_outputs = []
    for i in range(0, len(rx_signal), N):
        segment = rx_signal[i:i + N]
        if len(segment) == N:
            y = np.sum(segment * ref_signal)
            correlator_outputs.append(y)
            detected_bits.append(1 if y >= 0 else 0)
    return np.array(detected_bits), np.array(correlator_outputs)

detected_corr, y_corr = correlator_receiver(rx, s1, N)
bit_errors_corr = np.sum(bits_plot != detected_corr)

print("Transmitted bits:", bits_plot)
print("Detected bits (Correlator):", detected_corr)
print("Bit errors (Correlator):", bit_errors_corr)

# ============================================================
# PART D - MATCHED FILTER RECEIVER
# ============================================================
h_mf = s1[::-1]
mf_output = lfilter(h_mf, 1, rx)

sample_indices = np.arange(N - 1, len(mf_output), N)
mf_samples = mf_output[sample_indices]
detected_mf = (mf_samples >= 0).astype(int)
bit_errors_mf = np.sum(bits_plot != detected_mf[:len(bits_plot)])

print("Detected bits (Matched Filter):", detected_mf[:len(bits_plot)])
print("Bit errors (Matched Filter):", bit_errors_mf)

same_decisions = np.array_equal(detected_corr,
 detected_mf[:len(detected_corr)])
print("Correlator and matched filter equal:", same_decisions)

# ============================================================
# PART E - BER SIMULATION
# ============================================================
bits = np.random.randint(0, 2, n_bits)
symbols = 2 * bits - 1
tx_baseband = np.repeat(symbols, N)
ref_baseband = np.ones(N)

BER_sim = []
BER_theory = []

for EbN0_dB in EbN0_dB_range:
    EbN0_lin = 10 ** (EbN0_dB / 10)
    sigma2 = N / (2 * EbN0_lin)
    sigma = np.sqrt(sigma2)

    tx_ch = lfilter(h, 1, tx_baseband)
    rx_bb = tx_ch + sigma * np.random.randn(len(tx_ch))

    detected_bits = np.zeros(n_bits, dtype=int)
    for i in range(n_bits):
        seg = rx_bb[i * N:(i + 1) * N]
        y = np.sum(seg * ref_baseband)
        detected_bits[i] = 1 if y >= 0 else 0

    errors = np.sum(bits != detected_bits)
    BER_sim.append(errors / n_bits)
    BER_theory.append(0.5 * erfc(np.sqrt(EbN0_lin)))

BER_sim = np.array(BER_sim)
BER_theory = np.array(BER_theory)

print("Eb/N0 dB values:", EbN0_dB_range)
print("Simulated BER:", BER_sim)
print("Theoretical BER:", BER_theory)

# ============================================================
# PART F - SHANNON CAPACITY
# ============================================================
B = 2 * Rb
SNR_dB = np.arange(0, 21, 2)
SNR_lin = 10 ** (SNR_dB / 10)
C = B * np.log2(1 + SNR_lin)

snr_mark = 10
snr_mark_lin = 10 ** (snr_mark / 10)
C_mark = B * np.log2(1 + snr_mark_lin)

EbN0_req = 9.6
EbN0_shannon = 1.59
gap = EbN0_req - EbN0_shannon

print("Capacity at 10 dB SNR:", C_mark)
print("Gap to Shannon limit:", gap)
\end{lstlisting}
